% Generated draft: requires mathematical and source review.
\documentclass[11pt]{article}
\usepackage[margin=25mm]{geometry}
\usepackage{fontspec}
\setmainfont{Latin Modern Roman}
\usepackage{amsmath,amssymb,mathtools,unicode-math}
\setmathfont{Latin Modern Math}
\usepackage{enumitem,needspace,xurl,hyperref}
\setlength{\parindent}{0pt}
\setlength{\parskip}{5pt}
\setlength{\emergencystretch}{3em}
\newcommand{\N}{\mathbb{N}}
\newcommand{\Z}{\mathbb{Z}}
\newcommand{\Q}{\mathbb{Q}}
\newcommand{\R}{\mathbb{R}}
\newcommand{\C}{\mathbb{C}}
\newcommand{\F}{\mathbb{F}}
\newcommand{\A}{\mathbb{A}}
\newcommand{\PP}{\mathbb{P}}
\newcommand{\E}{\mathbb{E}}
\newcommand{\eps}{\varepsilon}
\newcommand{\dd}{\,\mathrm{d}}
\begin{document}
\section*{UnsolvedMath ID: 6800011}
\textbf{Source problem code:} catalogue entry
\subsection{Definitions and mathematical statement}
Let $M$ be a connected closed smooth manifold. The source asks whether
\[
\bigl[\exists\text{ Riemannian }g\text{ on }M:\ \operatorname{Sc}_g(x)>0
\text{ for every }x\in M\bigr]
\quad\Longrightarrow\quad
\bigl[M\text{ is }p\text{-atoral for every odd }p\bigr].
\]
No spin, orientability, fundamental-group, or dimension restriction is stated.

The record does not define $p$-atorality. A common mod-$p$ torality convention for an oriented closed $n$-manifold tests products of degree-one cohomology classes: torality would mean existence of $a_1,\ldots,a_n\in H^1(M;\F_p)$ with
\[
\langle a_1\smile\cdots\smile a_n,[M]_{\F_p}\rangle\ne0,
\]
and atorality would mean all such evaluations vanish. This is an indicated possible interpretation only; the source may involve covers or a different homological definition. For odd primes the displayed fundamental class requires orientability or a specified local-coefficient convention, which is missing. The phrase odd $p$ also does not explicitly say prime rather than another allowed modulus.

The supplied summary's language about essential tori or degree maps does not give a definition. Maps from a torus, maps to a torus, and mod-$p$ degree-one cup-product tests are different formulations. None is substituted for the source's undefined predicate without the original convention.
\subsection{Short English statement}
Does existence of positive scalar curvature force $p$-atorality for every odd parameter $p$? The meaning of atorality, prime/modulus range, and orientation conventions must be recovered from the source.
\subsection{Sources}
\begin{itemize}
\item Local UnsolvedMath registry, permanent ID 6800011, exact source code catalogue entry, ``Toral manifolds and positive scalar curvature---Question 2''; proposed by Bernhard Hanke, Augsburg, 2018. Registry status open, classification OPEN-TRIAGE.
\item Morgan and Pansu, \emph{A List of Open Problems in Differential Geometry} (2018), question environment 11, \url{https://www.imo.universite-paris-saclay.fr/~pansu/problems_MTDG.tex}. Supplied source-TeX extraction accessed July 29, 2026. The environment number and title question number are distinct.
\item The supplied summary mentions Hanke--Schick, index theory and Schoen--Yau-type arguments as related structural progress, without precise publications, definitions or hypotheses. It reports no definitive solution located. Its informal essential-torus/degree-map description is not treated as the definition of $p$-atorality.
\item Provenance flags current status and redistribution terms NEEDS\_REVIEW; source TeX is publicly available on a coauthor's publication page.
\end{itemize}
\end{document}
